\chapter{Packaging, deployment and observable services}
\label{evod44:chapter}

A model artifact is not yet a service. Suppose revision r1 serves requests
correctly and r2 changes a tokenizer option. The tensor dimensions still fit;
the readiness endpoint returns success; a kernel test passes. Should r2 accept
r1's checkpoints? Should it receive all traffic? Neither follows from those
observations. We need to identify the complete computation, define compatible
state and gather evidence about the service under a specified workload.

Our running example is a small release controller with two revisions, a bounded
number of admitted requests and synthetic completion times. It never starts a
service. It makes acceptance, promotion, draining and rollback executable
contracts. A real deployment must implement those contracts across crashes,
concurrency, authenticated artifact storage and device completion fences.

\section{Bind the release before starting it}
\index{release bundle}\index{canonical encoding}\index{provenance}
\label{evod44:bundle}
Define a bundle $B=(M,R,C,S,P)$: model identity $M$, runtime/compiler
identity $R$, execution configuration $C$, interpretation contracts $S$ and
provenance $P$. Tokenizer identity belongs beside model identity even if it
is stored in another file. $S$ includes API, state/checkpoint, IR and policy
schemas plus required backend features. The runtime is not free to reinterpret
an old native vector or KV tensor merely because its shape fits.

The reference serializes a fixed record in sorted-key JSON and computes
\[
d_B=\operatorname{SHA256}(\operatorname{encode}(B)).
\]
This encoding has no
floating-point or unordered-set ambiguity within its declared Python JSON
profile; it is not a claim of cross-language RFC 8785 canonicalization.
Production bundles should bind
digests of actual files, a canonical manifest format and dependency closure;
the fixture's strings are deliberately small substitutes. A mutable tag such
as ``latest'' can locate a manifest, but the resolved digest must identify
the revision used by each admitted request.

\Plate{01-artifact}{The manifest binds interpretation as well as payload. Changing a tokenizer or state schema creates a different identity even when weights are unchanged. Dashed arrows are verification inputs, not execution data.}{evod44:artifact}
\Listing{Bundle.digest}

SLSA describes provenance as verifiable information about an artifact's
production~\cite{slsa44}. Record source revision, builder, dependencies,
parameters and output digests. A nonempty provenance URI is not verification:
an attacker can type one. Verification needs a trusted issuer, subject-digest
binding and policy about acceptable builders. Our controller accepts a
\texttt{trusted\_provenance} boolean from such a hypothetical verifier;
it implements no signature verification and claims no SLSA level.

\subsection{A reproducibility test with two possible answers}
Build the same manifest twice: its digest must match. Change only the
tokenizer identifier: the digest must change. These are independent tests of
serialization stability and semantic closure. Neither proves that recompiling
source produces byte-identical binaries. Reproducible builds additionally
control toolchains, dependency resolution, timestamps and environment inputs.
Nor does artifact lineage prove model quality. Execution reproducibility,
build reproducibility and task validity answer different questions.

\subsection{Verify the delivered bytes against an independent identity}
\index{artifact integrity}\index{trusted pin}
A release capsule must deliver the payloads named by its manifest. The local
byte fixture has exactly four roles: model, runtime, tokenizer and configuration.
Each manifest field stores the SHA-256 of its corresponding immutable byte
string. The tokenizer payload is the four bytes \texttt{ACGT}; the other
payloads are small assigned strings, not executable binaries or learned weights.
The verifier rejects missing or extra roles, mutable payload containers, sizes
above one million bytes per role and any content-digest mismatch.
\Listing{verify_capsule}

First compare the manifest's identity with a trusted digest supplied through
an independent channel. Then check each payload against that manifest. These
checks answer different adversarial cases. Altering only weights fails the
payload check. Altering weights and updating their manifest digest is internally
consistent but changes the release identity, so the original trusted pin
rejects it. If the attacker can also supply a replacement trusted pin, both
checks pass: byte consistency cannot authenticate its own trust anchor. This
negative control prevents ``all hashes match'' from being mistaken for provenance.

The contract establishes exact delivered byte identity under collision resistance
and a trusted manifest pin. It does not parse model formats, reject NaN weights,
prove reproducible compilation or establish a model's task quality. Those are
separate gates. Chapter 28 already distinguishes tensor byte interpretation and
finite numerical validity. Here the capsule extends that principle to the
complete release envelope. The controller below still accepts metadata only:
an integration must run payload verification and backend loading before marking
the corresponding revision ready. No archive extraction or executable object
deserialization occurs in this data-only fixture.

\section{Compatibility is an acceptance relation}
\index{compatibility}\index{state schema}
\label{evod44:compat}
Write $\operatorname{Accept}(E,B)$ for an environment $E$ accepting bundle $B$.
It is a conjunction of required features and interpretation agreements, not
a comparison of version strings. In the reference it requires exact API,
state-schema, policy-schema and backend matches. This conservative rule
rejects potentially convertible states. Supporting a converter requires a
versioned transformation, old/new fixtures and a rollback policy.

For Chapter 31's restricted diagonal recurrence, bind the embedding, retention,
write and readout parameters, dtype, vocabulary and consumed position. For
Chapter 35's actual Chapter 23 Hermon stack, bind learned absolute positions,
per-layer K/V geometry, normalization and causal mode. Replacing that stack
with a rotary single-head demonstration would change the artifact. DOGMA
remains the non-Transformer direction; Hermon DNA remains Transformer-based.
This controller inspects release metadata; it does not verify either engine's
numerical state or integrate an unfinished Chapter 39 runtime.

\section{Roll out a revision, not a filename}
\index{canary}\index{shadow execution}\index{rollback}
\label{evod44:rollout}
A registered revision has passed metadata gates. It is not ready until its
weights and dependencies are usable. A canary is eligible for a declared small
traffic cohort; active means selected for normal traffic. Superseded revisions
stop receiving new work but may finish old requests. Draining closes admission
before termination. A rolled-back candidate likewise retains responsibility
for work already admitted.

\Plate{02-rollout}{Promotion changes the routing choice. Existing request ownership remains attached to its admitted revision. Rollback closes the candidate's admission without erasing its unfinished requests.}{evod44:rolloutplate}

Canary percentage, workload strata and observation window must be chosen
before looking at results. A shadow candidate can run without publishing its
output, but it must also be prevented from duplicating external effects.
Shadow traffic that omits expensive tools or uses only short prompts cannot
certify the corresponding full service workload.

\Listing{Deployment.promote}
The supplied package used global telemetry for promotion and allowed an empty
sample to pass. The revised gate uses only the candidate's own completed
observations and requires a minimum count. This prevents r1's healthy history
from certifying an unobserved r2. It is still a simple serial specification:
a real canary needs windows, representative traffic and unresolved-request
deadlines, since counting only completed requests can hide stuck ones.

\section{Live, ready and able to reserve are separate predicates}
\index{liveness}\index{Readiness}\index{admission}\index{draining}
\label{evod44:admit}
Liveness asks whether a process can continue; readiness asks whether it should
receive traffic. Startup may include loading or compilation. Kubernetes
distinguishes these probe purposes~\cite{probes44}; our service contract uses
that distinction without implementing Kubernetes. A live process can be unready
while warming, losing a dependency, exhausting memory or draining. Overload
should not automatically trigger a liveness restart that discards useful work.

Even a ready revision can reject a request because capacity is exhausted.
For request $r$, admission requires readiness, a compatible class, tenant
authority, capacity and a plausible deadline. The simple controller reserves
one global slot; the later isolation contract reserves resource vectors.
Real estimates include prompt work, output reservation, native state or KV,
tool budget and downstream decode capacity.

\Listing{Deployment.admit}
\Listing{Deployment.complete}

Completion validates the observation before changing ownership, and binds an
issued request ID to the revision that admitted it. Duplicate, foreign and
late duplicate completions fail. This makes accounting testable instead of
trusting callers to decrement a counter exactly once.

\subsection{The conservation law behind graceful draining}
Let $A_d$ be cumulative admitted requests, $C_d$ completed requests and $I_d$
currently owned requests for revision $d$. All are integer counts, not rates.
Initially they are zero. Admission increments $A_d,I_d$; completion increments
$C_d$ and decrements $I_d$. No other operation changes these quantities. Thus
\begin{equation} A_d=C_d+I_d. \label{evod44:conserve}\end{equation}
This is an induction over accepted events. Rejected operations preserve all
three values. Closing admission prevents further increments of $A_d$;
termination requires $I_d=0$. A timeout is not an acknowledgement that a GPU
kernel or tool no longer owns memory. A production service must add explicit
cancel/acknowledge and output-publication accounting from Chapter 30.

\Plate{03-drain}{Drain is a protocol over ownership. At the middle instant admission is closed, but two issued requests still prevent termination. Late valid completions reduce ownership; duplicate completions cannot do so twice.}{evod44:drain}

\section{Observe causal operations without copying private content}
\index{observability}\index{trace context}\index{telemetry cardinality}
\label{evod44:observe}
A request may cross gateway, planner, compiler, router, engine, memory and
tools. Parent/child spans represent causal operations; their durations can
overlap, so summing them can exceed end-to-end latency. OpenTelemetry supplies
common trace, metric and log concepts~\cite{otel44}. Domain attributes should
bind artifact, logical engine, phase and outcome using versioned names.
Trace context is correlation data, not authentication or tenant authority.

Track request/error rates, duration distributions, queue delay, TTFT, ITL and
useful tokens/s alongside admission rejection. Record batch occupancy,
KV/native-pool utilization, cache hit rates, retrieval/tool latency and rollback
events so an SLO change can be related to resource pressure and plan changes.
These are instrumentation requirements; the fixture emits only its small
lifecycle schema and assigned completions, not these live metrics.

Use bounded-cardinality metric labels for engine, revision class and outcome.
Putting request IDs, prompts or arbitrary resource paths into metric labels
can exhaust the telemetry backend. Keep a restricted correlation ID in traces
and diagnostic logs. Sampled traces cannot alone count every failed request;
use unsampled counters for SLI denominators and record dropped telemetry.
Redaction must happen before export, including exception messages and span
attributes. Prompts, outputs, retrieved private documents and credentials
are absent from our event schema.

The reference's sequential event numbers describe one process. They are not
tamper-evident storage or distributed ordering. A production audit sink needs
access control, retention, delivery/failure policy and privacy review.

\section{Turn observations into a falsifiable service objective}
\index{SLO}\index{percentile}\index{Error budget}
\label{evod44:slo}
An SLI is an observed indicator; an SLO is its target over an explicit
population and window. For eligible offered requests in a one-hour class,
an example is $P(\text{success}\land TTFT\le500\,\mathrm{ms})\ge0.99$.
State how rejections, timeouts, cancellations and maintenance count.
Accepted-only latency can improve merely by rejecting difficult requests.

The reference has a simpler completion-cohort gate: error fraction at most
0.01 and nearest-rank p95 latency at most 500 ms after at least 20 observations.
For sorted values $x_{(1)}\le\cdots\le x_{(n)}$, the estimator is
$x_{(\lceil qn\rceil)}$, $0<q\le1$. Three durations $(10,20,900)$ have
p95 900 ms. A lower interpolation/index shortcut could conceal that tail.
\Listing{Gate.passes}
\begin{center}\input{\Generated/gates.tex}\end{center}
The generated teaching trace deliberately overrides the minimum to three:
both r1 and r2 observe $(40,60,130)$ ms, then r2 is rolled back. These are
assigned numbers, not a reliability experiment or a recommended sample size.

For target $A^*$, allowed bad fraction is $E=1-A^*$. With $N$ eligible
requests, the budget is $NE$ bad requests; burn rate is
$b=(N_{bad}/N)/E$. At $A^*=0.99$ and $N=10{,}000$, 100 bad requests consume
the budget; 200 imply $b=2$. A short window estimates burn noisily. Even zero
failures among $n$ independent Bernoulli trials has one-sided 95\% upper
failure bound $1-0.05^{1/n}$: about 0.139 for $n=20$. Independence and a
stationary representative population are assumptions, often violated by
correlated outages. Passing the deterministic gate is no confidence guarantee.

\section{Configuration, secrets and a rollback that can actually run}
\label{evod44:rollback}
Execution-critical configuration is immutable in the bundle. Dynamic policy
knobs need version, validation and an effective-change event; admission must
bind the applicable version. A secret reference can be configuration, but its
resolved value must remain outside bundles, prompts, provenance and ordinary
telemetry. Rotation needs scoped runtime resolution, not rebuilt model weights.
This chapter resolves no secret and contacts no service.

Rollback changes executable selection. It does not reverse emails, payments,
deletions or other effects already committed by tools. Retain old immutable
artifacts and verify that the current durable-state schema is readable by
the target. If r2 irreversibly migrated durable memory, reverting binaries
alone can make recovery worse. A tested reversible migration, forward-compatible
reader or separate restoration plan is required.
\Listing{Deployment.rollback}

The reference refuses an incompatible or unavailable target, closes candidate
admission and preserves its issued requests. It models one rollback pair,
not a fleet consensus protocol. Crashes between routing and ownership updates,
split controllers, lost journals and device completion races remain outside it.

\section{Break the contract deliberately}
\label{evod44:failure}
Tests change semantic manifest fields, reject false provenance verification,
attempt empty canary promotion, contaminate one revision with another's
observations, exhaust global capacity, duplicate completions, drain with work
outstanding and roll back into incompatible state. A bounded enumeration
checks all $4^5=1024$ five-event words over admit/complete/drain/terminate
against an independently maintained owner set. It catches accounting and
state-machine errors, not arbitrary concurrent interleavings or process death.

Production experiments should additionally interrupt artifact download, kill
the controller at each journal boundary, lose the telemetry sink, race shutdown
with kernel acknowledgement and replay a routing update. Define expected
recovery before injecting each fault. No such live experiments are run here.

\section{Twelve worked exercises}
\begin{enumerate}
\item \textbf{Bind configuration.} Weights are unchanged, but tokenization changes.
May the digest remain equal? \emph{Solution:} no; the computed token sequence
and state interpretation can change. Include tokenizer/config identity.
\item \textbf{Separate provenance and quality.} A trusted builder produced poor
weights. What passes? \emph{Solution:} lineage verification may pass while task
evaluation fails; require both gates rather than treating either as the other.
\item \textbf{Predict readiness.} Weights are loading in a healthy process.
\emph{Solution:} live, unready and closed to admission; liveness alone is insufficient.
\item \textbf{Drain arithmetic.} $A=11,C=8$. May shutdown finish?
\emph{Solution:} $I=3$ by Equation~\ref{evod44:conserve}; no. First close admission,
then settle all three ownership obligations.
\item \textbf{Correlate logs.} Why include a safe trace ID but omit prompt text?
\emph{Solution:} the ID links causal diagnostics without copying protected payloads.
It still needs access controls and is not an authorization credential.
\item \textbf{Calculate the tail.} Find nearest-rank p95 of $(10,20,900)$.
\emph{Solution:} rank $\lceil2.85\rceil=3$, hence 900 ms; the 310 ms mean hides it.
\item \textbf{Test the SLO claim.} A kernel is twice as fast. Is the service SLO proved?
\emph{Solution:} no; queue, transfer, tools and offered-load outcomes remain unmeasured.
\item \textbf{Undo an effect.} r2 sent an email before rollback.
\emph{Solution:} selecting r1 does not retract it. Use an effect-specific compensation
or idempotency protocol; never promise generic undo.
\item \textbf{Check schemas.} Old code reads schema 1; durable memory is schema 2.
\emph{Solution:} fail rollback unless a tested compatible reader/restoration exists.
\item \textbf{Place secrets.} Should a build manifest include an API token?
\emph{Solution:} no; include a scoped reference when needed and resolve after authorization.
\item \textbf{Quantify canary evidence.} Zero failures in 20 independent trials proves
99\% reliability? \emph{Solution:} no; the stated upper failure bound is about 13.9\%.
Representativeness and dependence are additional problems.
\item \textbf{Design a release experiment.} \emph{Rubric:} freeze bundle and workload;
declare compatibility, cohort, minimum evidence and timeout accounting; verify
readiness/admission/drain; redact correlated telemetry; test rollback and durable
schema; archive raw outcomes and rejected hypotheses. A checklist without outcomes
earns no claim of a successful deployment.
\end{enumerate}

\section{Terms and the next boundary}
\Term{Release identity} binds the execution-critical manifest.
\Term{Readiness} permits routing but does not reserve capacity.
\Term{Error budget} is allowed bad service over a declared population/window.
Chapter 45 adds mutually untrusted tenants. The same ownership and telemetry
must now preserve identity, authority, resource quotas and privacy.
